% !TEX root = ../Thesis.tex

\chapter{Implementazione del sistema}
\label{ch:implementazione}

Il presente capitolo descrive l'implementazione pratica del sistema steganografico presentato nel \autoref{ch:proposta}. Mentre il capitolo precedente si è concentrato sugli aspetti algoritmici e teorici — in particolare gli \autoref{alg:selezione-token}, \ref{alg:codifica-patch} e \ref{alg:codifica-messaggio} per l'encoding steganografico, e gli \autoref{alg:decodifica-token} e \ref{alg:decodifica-messaggio} per il decoding steganografico — qui vengono illustrati i dettagli realizzativi: l'architettura software, l'organizzazione del codice in classi Python, e le procedure per la generazione e il salvataggio delle immagini.

\section{Architettura software}
\label{sec:architettura-software}

\subsection{Scelta del linguaggio di programmazione e del framework}

L'intero sistema è stato sviluppato in Python con il framework \emph{PyTorch}. Sia l'implementazione originale di VQ-GAN \cite{esser-2020} sia quella di \emph{Meteor}~\cite{kaptchuk-2021} sono sviluppate in questo ecosistema, la cui adozione consente di riutilizzare direttamente i pesi pre-addestrati dei modelli esistenti e di garantire la riproducibilità dei risultati. Per motivi strettamente pratici non sono state prese in considerazione alternative come TensorFlow o JAX, che avrebbero richiesto una riscrittura completa del codice e una conversione dei modelli, con conseguente rischio di introdurre discrepanze nei risultati, senza una relae giustificazione in termini di prestazioni o funzionalità specifiche.

\subsection{Moduli principali}

Il sistema implementato si compone di tre moduli principali, ciascuno corrispondente a un file sorgente:

\begin{description}
	\item[\texttt{vqgan.py}:] interfaccia con il modello VQ-GAN pre-addestrato. Espone le funzioni \texttt{encode\_to\_z()} e \texttt{decode\_to\_img()}, che astraggono rispettivamente l'encoding VQ-GAN di un'immagine nella corrispondente griglia di indici discreti e il decoding VQ-GAN di una griglia di indici nell'immagine RGB finale.
	
	\item[\texttt{encode\_methods.py}:] implementa la classe \texttt{SteganographyEncoder}, che realizza la fase di encoding steganografico descritta nell'\autoref{sec:encoding-algorithm}. Gestisce l'inizializzazione dei tensori, l'iterazione sulle posizioni della griglia, la selezione dei token e la verifica di codifica.
	
	\item[\texttt{decode\_methods.py}:] implementa la classe \texttt{SteganographyDecoder}, che realizza la fase di decoding steganografico descritta nell'\autoref{sec:decoding-algorithm}. Ricostruisce le distribuzioni di probabilità percorrendo la griglia e recupera i bit codificati.
\end{description}

\section{Le classi Steganography~Encoder e Steganography~Decoder}
\label{sec:classi-encoder-decoder}

\subsection{SteganographyEncoder}

La classe \texttt{SteganographyEncoder} coordina la generazione guidata dal messaggio. Il suo metodo principale \texttt{encode\_message\_to\_image()} implementa l'\autoref{alg:codifica-messaggio} nei seguenti passi:

\begin{enumerate}
	\item Estrae l'immagine di contesto dal dataset condiviso e la codifica tramite l'encoder VQ-GAN (\texttt{encode\_to\_z()}) per ottenere la matrice di riferimento $\mathbf{R}$.
	\item Inizializza il tensore in costruzione $\mathbf{B}$ copiando le prime $R$ righe da $\mathbf{R}$ e azzerando le restanti posizioni.
	\item Converte il messaggio in sequenza binaria $\mathbf{b}$.
	\item Itera sulle posizioni $(r,c)$ della griglia in ordine raster: per ciascuna, costruisce il contesto $\mathbf{c}$ come descritto nella \autoref{sec:gestione-contesto}, invoca la selezione steganografica del token (\autoref{alg:selezione-token}) ed esegue la verifica di codifica (\autoref{alg:codifica-patch}).
	\item Quando $\mathbf{b}$ è esaurito, completa le posizioni rimanenti con campionamento casuale per garantire la coerenza visiva dell'immagine.
	\item Decodifica $\mathbf{B}$ tramite il decoder VQ-GAN (\texttt{decode\_to\_img()}) per ottenere l'immagine finale e la salva su disco.
\end{enumerate}

\subsection{SteganographyDecoder}

La classe \texttt{SteganographyDecoder} implementa l'\autoref{alg:decodifica-messaggio}. Il suo metodo \texttt{decode\_message()} percorre la griglia nello stesso ordine raster dell'encoding steganografico: per ciascuna posizione $(r,c)$, ricostruisce il contesto $\mathbf{c}$, estrae il token $z^*$ dall'immagine ricevuta tramite l'encoder VQ-GAN (\texttt{encode\_to\_z()}), e invoca la decodifica steganografica del token (\autoref{alg:decodifica-token}) per recuperare i bit. I bit così ottenuti vengono concatenati e riconvertiti nel messaggio originale.

\section{Costruzione del contesto}
\label{sec:costruzione-contesto}

Come illustrato nella \autoref{sec:gestione-contesto}, il Transformer opera su una finestra locale $p \times p$. L'implementazione di questa costruzione si articola in due funzioni ausiliarie. \texttt{local\_indexes()} calcola le coordinate della finestra $p \times p$ centrata sulla posizione corrente, gestendo i casi di bordo mediante traslazione della finestra per mantenerla interamente entro i limiti della griglia. \texttt{build\_context\_from\_patches()} utilizza tali coordinate per estrarre le porzioni corrispondenti dalla matrice di riferimento $\mathbf{R}$ e dal tensore in costruzione $\mathbf{B}$, concatenandole in un'unica sequenza di $2 \times p \times p$ token da fornire al Transformer.

\section{Generazione e salvataggio delle immagini}
\label{sec:generazione-immagini}

\subsection{Verifica di codifica}

Un aspetto implementativo critico è il meccanismo di verifica descritto nell'\autoref{alg:codifica-patch}. L'implementazione esegue un ciclo di tentativi in cui per ogni token selezionato $i^*$:
\begin{enumerate}
	\item Inserisce $i^*$ nel tensore $\mathbf{B}$ alla posizione corrente;
	\item Decodifica $\mathbf{B}$ in un'immagine parziale tramite il decoder VQ-GAN (\texttt{decode\_to\_img()});
	\item Ricodifica l'immagine parziale tramite l'encoder VQ-GAN \\(\texttt{encode\_to\_z()});
	\item Confronta l'indice ricostruito $\hat{\imath}$ con $i^*$.
\end{enumerate}
Se $\hat{\imath} = i^*$ la codifica è considerata valida; altrimenti si ripete con un nuovo tentativo. Dopo 100 tentativi falliti, la selezione viene forzata per evitare un ciclo infinito, accettando il token corrente anche in caso di discordanza. Questo costo computazionale, come discusso nella \autoref{sec:limitazioni}, rende l'encoding steganografico significativamente più oneroso del decoding steganografico.

\subsection{Salvataggio in formato PNG}

Le immagini generate vengono salvate in formato PNG (Portable Network Graphics). Tale scelta non è accidentale: il formato PNG è \emph{lossless} e supporta una profondità di colore di 8 bit per canale (24 bit totali per pixel in RGB), preservando integralmente ogni valore dei tre canali di colore senza perdita di informazione. Ciò è essenziale per la corretta decodifica steganografica del messaggio, poiché qualsiasi compressione con perdita — come quella del formato JPEG, che opera una quantizzazione delle componenti di frequenza nello spazio YCbCr — altererebbe i valori dei pixel e, di conseguenza, gli indici discreti ottenuti dalla ricodifica tramite l'encoder VQ-GAN (\texttt{encode\_to\_z()}), compromettendo il recupero del messaggio. Il formato PNG garantisce invece che l'immagine generata possa essere salvata, trasmessa e ricaricata con degradazione minimale. Non è stata effettuata un'analisi empirica della resitenza del sistema a formati con perdita. un tale studio potrebbe essere oggetto di ricerche future.

\section{Limitazioni pratiche}
\label{sec:limitazioni-pratiche}

Oltre alle limitazioni teoriche discusse nella \autoref{sec:limitazioni}, l'implementazione presenta alcune criticità.

\subsection{Costi computazionali}

Il ciclo di verifica della codifica rende il processo di encoding steganografico computazionalmente oneroso: ogni token può richiedere fino a 100 cicli di decoding VQ-GAN e successivo encoding VQ-GAN dell'immagine parziale. Per un'immagine $256 \times 256$ con circa $240$ token da codificare, il numero di operazioni coinvolte può diventare significativo, con un rapporto asimmetrico tra i tempi di encoding steganografico e decoding steganografico.

\subsection{Gestione dei casi di bordo}

La traslazione della finestra di contesto vicino ai bordi della griglia, pur garantendo che la finestra sia sempre di dimensione $p \times p$, sacrifica la simmetria centrata sulla posizione corrente per le celle prossime al bordo.

\subsection{Riproducibilità dei risultati}

L'uso di un seme pseudo-casuale fisso garantisce la riproducibilità del campionamento casuale nelle posizioni non steganografiche. Tuttavia, operazioni di arrotondamento e quantizzazione nell'encoder VQ-GAN possono introdurre variazioni minime tra esecuzioni su hardware differente (diverse implementazioni GPU di operazioni floating-point).